% Lösungen Einheit 2 von 4 – Parameter aus Lagebedingungen (zusammenbau v0.1)
\einheitenkopf{Einheit 2 von 4 · Parameter aus Lagebedingungen}

\erg{12}{a) Normalenvektor – die Stellung der Ebene ändert sich \quad b) $\vec{n}_k \cdot \vec{u} = 2k + 2 - 4 = 0$, also $k = 1$ \quad c) $(4 | k | -2) = 2 \cdot (2 | 1 | -1)$, also $k = 2$ \quad d) ja: Richtung $k - 1 = 0$ gibt $k = 1$; Stützpunkt $k + 2 + 1 = 4$ gibt $k = 1$ – beide passen, $g$ liegt in $E_1$ \quad e) $P$: $2 + 2 + 3 = 7$, $Q$: $4 + 0 + 3 = 7$ – beide liefern $k = 7$, die Kante liegt in $E_7$ \quad f) $M(0 | 5 | 0)$; aus $x_1$ und $x_3$ folgt $s = 2$, aus $x_2$ dann $2k - 3 + 2 = 5$, also $k = 3$ \quad g) $E: 4x_1 + 3x_2 + 2x_3 = 11$; $4k + 9 + 4 - 2k = 11$, also $k = -1$ \quad h) Benachbart hieße: $\overrightarrow{QR}$ parallel zu $(2 | 1 | 2)$ (dieselbe Gerade) oder senkrecht dazu (zwei Geraden). $\overrightarrow{QR} = (3 | 3 | 0)$ ist kein Vielfaches, und $(3 | 3 | 0) \cdot (2 | 1 | 2) = 9 \ne 0$ – beide Fälle scheiden aus}
\erg{13}{a) Fehler: der Stützpunkt wurde in die Ebene eingesetzt. Richtig: $\vec{n}_k \cdot \vec{u} = 2k + 1 + k - 3 = 0$, also $k = \tfrac{2}{3}$ \quad b) Der Normalenvektor steht senkrecht auf der Ebene: Ist $g$ parallel zur Ebene, steht ihr Richtungsvektor senkrecht auf dem Normalenvektor; steht $g$ senkrecht auf der Ebene, zeigt ihr Richtungsvektor in Richtung des Normalenvektors.}
